\documentclass[11pt]{article}

\usepackage{geometry}
\geometry{margin=1in}

\usepackage{amsmath,amssymb,amsfonts,amsthm}
\usepackage{mathtools}
\usepackage{hyperref}
\usepackage{graphicx}
\usepackage{float}
\usepackage{booktabs}
\usepackage{tikz-cd}
\usepackage{cleveref}
\usepackage{subcaption}
\usepackage{placeins}
\usepackage{enumitem}
\usepackage{array}
% xr provides \externaldocument so cross-document \ref{S-...} works.
% Plain xr (not xr-hyper) is used because xr-hyper imports labels in a
% hyperref-extended format that cleveref's templabel parser does not accept.
% The cost is that cross-document references show as text only rather than
% as live hyperlinks; the in-document \ref number resolution still works
% and there are no undefined-reference errors.
\usepackage{xr}

\hypersetup{colorlinks=true, linkcolor=blue, citecolor=blue, urlcolor=blue}

%---------------------------
% Theorem environments
%---------------------------
\numberwithin{equation}{section}

\newtheorem{theorem}{Theorem}[section]
\newtheorem{lemma}[theorem]{Lemma}
\newtheorem{proposition}[theorem]{Proposition}
\newtheorem{corollary}[theorem]{Corollary}
\theoremstyle{definition}
\newtheorem{definition}[theorem]{Definition}
\newtheorem{assumption}[theorem]{Assumption}
\theoremstyle{remark}
\newtheorem{remark}[theorem]{Remark}

\newcommand{\proofsuppref}[1]{%
  \par\noindent\emph{Proof.} See the online supplement, Section~\ref{S-#1}.\par\medskip
}
\newcommand{\proofappendixref}[1]{\proofsuppref{#1}}


%---------------------------
% Notation shortcuts
%---------------------------
\newcommand{\K}{\mathcal{K}}
\newcommand{\Hil}{\mathcal{H}}
\newcommand{\bd}{\partial}
\newcommand{\Id}{\mathbf{1}}
\newcommand{\Proj}{\mathbf{P}}
\newcommand{\spec}{\mathrm{spec}}
\newcommand{\dist}{\mathrm{dist}}
\newcommand{\rank}{\mathrm{rank}}
\newcommand{\tr}{\mathrm{tr}}
\newcommand{\Ran}{\mathrm{Ran}}
\newcommand{\opnorm}[1]{\left\lVert#1\right\rVert}
\newcommand{\e}{\varepsilon}
\newcommand{\sgn}{\mathrm{sgn}}
\newcommand{\Indint}{\operatorname{Ind}_{\mathrm{int}}}
\newcommand{\Ch}{\operatorname{Ch}}
\newcommand{\hexagon}{\mathrm{hex}}

% Resolve cross-document \ref{S-...} labels via xr.
\externaldocument[S-]{paper_supplement}

%---------------------------
% Example name macros
%---------------------------
\newcommand{\ExOneName}{Disk-like patch (1 boundary)}
\newcommand{\ExTwoName}{Annulus-like band (2 boundaries)}
\newcommand{\ExThreeName}{Three-boundary patch (disk with two holes)}

%---------------------------
% Figure safety helper
%---------------------------
\newcommand{\safeincludegraphics}[2][]{%
  \IfFileExists{#2}{\includegraphics[#1]{#2}}{%
    \fbox{\begin{minipage}[c][0.18\textheight][c]{0.92\linewidth}
      \centering Missing figure file\\[0.3em]\texttt{\detokenize{#2}}
    \end{minipage}}%
  }%
}

\title{Boundary and Interface Classes for Universal Chern Hamiltonians\\
via Doubling and Jamming}

\author{
Yitzchak Shmalo\thanks{Corresponding author: \texttt{yitzchak.shmalo@gmail.com}.}\\[0.6em]
Einstein Institute of Mathematics\\
Hebrew University of Jerusalem
}
\date{\today}

\begin{document}
\maketitle

\begin{abstract}
Higson and Prodan's universal Chern Hamiltonians are closed-surface models built from the
simplicial boundary operator and Poincar\'e duality. We treat the boundary and interface classes
obtained from finite surfaces with boundary by doubling the surface and jamming the reflected copy.
For every finite triangulated oriented surface with boundary we construct an intrinsic
Poincar\'e--Lefschetz Hamiltonian, compare it with the original-side compression of the doubled
closed Hamiltonian, identify the discrepancy as boundary-collar supported, and prove
Schur-complement estimates for finite-\(M\) jamming. Along families with negligible collar
dimension these estimates identify the normalized bulk spectral limits of the intrinsic,
compressed, and jammed models. The headline Roe-algebra theorem identifies the doubled-and-jammed
interface class in \(K_1(C^*(Y\subset X))\) with
\(\partial_Y(E_+(\pi(P_\pm)))\), using the Higson--Prodan closed universal-model theorem
\cite{HigsonProdanUniversalPRL,HigsonProdanUniversalArxiv} for the closed-bulk gap and Chern
value; this \(K_1\)-class identity is unconditional modulo HP. The integer interface winding
\(\pm1\) is proved internally in the periodic triangular Toeplitz case via the
\(K_1\)-to-\(\mathbb Z\) Toeplitz extension, and is stated as a conditional corollary for
general partitions, conditional on the existence of an integer-valued Chern-jump pairing for
the interface ideal. The periodic triangular Bloch class is a fully internal benchmark: its
closed-bulk gap, Chern signs, and straight Toeplitz winding are proved in the paper. Companion
numerical code and the online supplement are available at the project repository
(\url{https://github.com/yspennstate/Doubling-and-Jamming-A-Universal-Chern-Hamiltonian-on-Triangulated-Surfaces-with-Boundary});
additional hexagonal and square two-band symbols, robustness variants, full appendix proofs,
and the numerical record are collected in that online supplement.
\end{abstract}

%====================================================================
\section{Introduction}
%====================================================================

\subsection{Universal Chern Hamiltonians from simplicial data}
The Haldane model is the prototypical two-dimensional class-A topological insulator
\cite{Haldane1988}; for background on the tenfold-way classification and Chern phases, see
\cite{SchnyderRyuFurusakiLudwig,KitaevPeriodic,QiZhangRMP,ProdanSchulzBaldesBook}. On irregular
discretizations, direct magnetic-field or Peierls constructions can lead to mobility-gap or
defect-state issues that complicate the extraction of clean topological gaps; see, for example,
\cite{Mitchell2018,BourneProdan2018,BourneMesland2019,ProdanSchulzBaldesBook}. Rigorous
bulk--edge and bulk--interface correspondences in quantum Hall and related class-A systems are
treated in
\cite{ElbauGraf2002,ElgartGrafSchenker2005,GrafPorta2013,BourneKellendonkRennie,KubotaControlled,EwertMeyer2019,AlldridgeMaxZirnbauer2020,LudewigThiangEdge}.

Higson and Prodan introduced a different route
\cite{HigsonProdanUniversalPRL,HigsonProdanUniversalArxiv}. For a triangulation of a \emph{closed}
oriented surface, one builds a local tight-binding Hamiltonian from two pieces of
simplicial-topological data: the simplicial boundary operator $B$ and a Hermitian Poincar\'e-duality
operator $S$ constructed from cap product with the fundamental cycle. The resulting Hamiltonians
\begin{equation}\label{eq:Hpm_intro}
H_\pm=(B+B^\dagger)\pm S
\end{equation}
are local and, in the sequential-refinement limit, define nontrivial topological bulk classes with
Chern numbers $\pm1$ \cite{HigsonProdanUniversalPRL,HigsonProdanUniversalArxiv}. The underlying
Hilbert--Poincar\'e framework uses the algebraic identity \(BS+SB^\dagger=0\), rooted in analytic
surgery theory and the work of Higson and Roe
\cite{HigsonRoeI,HigsonRoeII,HigsonRoeIII,WallSurgery}.

\subsection{The boundary case}
If the input triangulation $\K$ has boundary, there is no closed fundamental $2$-cycle supported
on the faces of $\K$, so the closed-case cap-product recipe for $S$ is not directly available.
Compact manifolds with boundary instead carry a relative fundamental class and satisfy
Poincar\'e--Lefschetz duality \cite{Hatcher}. We use this relative class to define an intrinsic
boundary Hamiltonian on $\K$, and then compare it with a second construction that reduces the
boundary problem to the closed case by simplicial doubling and jamming the reflected side.

\subsection{Doubling and jamming}
Given a triangulated surface-with-boundary $\K$, let $\overline\K$ be an orientation-reversed copy.
Gluing $\K$ and $\overline\K$ along $\bd\K$ produces the closed double $D(\K)$. We then build the
closed universal Hamiltonians on $D(\K)$ and suppress the reflected copy by adding a large on-site
potential on the copy degrees of freedom. This is the \emph{doubling-and-jamming} construction. The
finite-\(M\) analysis gives error bounds for this approximation. For refinement families whose
closed doubles are covered by Higson--Prodan's published closed theorem, that theorem supplies the
closed-bulk gap and Chern class. Section~\ref{sec:pathb_internalization} proves, for the
original/copy partition, that the Roe connecting map used to define the interface class is the
Ewert--Meyer bulk--interface map for the doubled-jammed data.

\subsection{Theorem map}
Let
\[
H_{\pm,M}=H_\pm(D(\K))+M\Proj_{\mathrm{copy}}
\]
be the jammed Hamiltonian, written in block form as
\[
H_{\pm,M}=
\begin{pmatrix}
A&C\\
C^\dagger&D+M\Id
\end{pmatrix}
\quad\text{on}\quad
\Hil_{\mathrm{orig}}\oplus\Hil_{\mathrm{copy}}.
\]
\begin{enumerate}[leftmargin=2em]
\item The relative fundamental class of \((\K,\bd\K)\) gives an intrinsic
Poincar\'e--Lefschetz boundary Hamiltonian on \(\K\). The simplicial double gives a second
Hamiltonian by closed Higson--Prodan theory; its original-side compression agrees with the
relative Hamiltonian up to a boundary-collar operator, and the jammed operator converges to this
compression at low energy.
\item The finite-\(M\) comparison is made quantitative by Schur-complement estimates:
compressed resolvents, full block resolvents, first-order effective operators, bounded-energy
spectral localization, leakage, Morse-index stability, and contour spectral projections. Combined
with the collar-rank comparison, these estimates give a fully internal bulk spectral-measure
equivalence between the intrinsic Poincar\'e--Lefschetz model, the doubled compression, and the
finite-\(M\) jammed model along any family with negligible boundary-collar dimension and
\(M\to\infty\).
\item For refinement families whose closed doubles are covered by the Higson--Prodan theorem, the
HP gap gives the quotient by the interface ideal and hence the global Roe-algebra interface class
\[
\Indint(J_{\pm,M}^\infty;T^+,T^-)
=
\partial_Y\!\bigl(E_+(\pi(P_\pm))\bigr)
\in K_1(C^*(Y\subset X)).
\]
Theorem~\ref{thm:pathb_roe_winding}, proved in Section~\ref{sec:pathb_internalization}, identifies
this class as \(\partial_Y(E_+(\pi(P_\pm)))\). In the straight periodic Toeplitz case the
integer pairing is internal (Corollary~\ref{cor:toeplitz_winding}) and gives winding
\(\pm1\); for general partitions the integer pairing is recorded as a conditional corollary
(Corollary~\ref{cor:conditional_winding}).
\item For the periodic triangular-lattice class, the closed-bulk gap and Chern input are checked
directly. We write down the \(6\times6\) Bloch symbol, prove a uniform zero gap from an exact
determinant formula, compute the Chern signs from the isolated massless crossing, and evaluate
primitive straight periodic interfaces by the Toeplitz extension. Finite covers, annular families,
finite-corner quotients, finite-defect variants, and tail-split component maps are treated in the
online supplement, Section~\ref{S-app:section5_variants}.
\end{enumerate}

\paragraph{Logical dependencies.}
The proof uses the following implications; no converse is assumed.
\[
\begin{array}{ll}
\mathrm{(I)} & \text{internal finite-dimensional results: the relative boundary operator,}\\
& \text{the collar comparison, and the large-\(M\) Schur-complement estimates;}\\[0.2em]
\mathrm{(HP)} & \text{the Higson--Prodan closed-bulk gap and Chern theorem for the doubles;}\\[0.2em]
\mathrm{(EM)} & \text{the internal identification of the original/copy Roe boundary map}\\
& \text{with the Ewert--Meyer pairing (Theorem~\ref{thm:pathb_roe_winding}).}
\end{array}
\]
From (I) and the gap part of (HP), the Roe construction supplies the global interface
\(K_1\)-class. From the Chern part of (HP) and (EM), Theorem~\ref{thm:pathb_roe_winding}
identifies the \(K_1\)-class as \(\partial_Y(E_+(\pi(P_\pm)))\). Thus the headline conclusion
of Theorem~\ref{thm:pathb_roe_winding} is the unconditional
\(K_1\)-class identity
\(\mathrm{(I)}+\mathrm{(HP)}+\mathrm{(EM)}\Longrightarrow
\Indint=\partial_Y(E_+(\pi(P_\pm)))\). The integer evaluation
\(\operatorname{Wind}=\pm1\) is internal in the periodic triangular Toeplitz case
(Corollary~\ref{cor:toeplitz_winding}) and conditional in general
(Corollary~\ref{cor:conditional_winding}). The ingredients (I) and (EM) are proved here. The
ingredient (HP) is exactly the cited closed theorem of Higson--Prodan. In the periodic
triangular section, the gap and Chern signs are proved directly, and primitive straight
periodic interfaces are evaluated internally by the Toeplitz extension.

\begin{table}[H]
\centering
\small
\caption{Theorem map. ``Internal'' means proved here from the inputs shown; ``modulo HP'' means
that the only cited closed-bulk theorem used is Higson--Prodan's published PRL/arXiv theorem.}
\label{tab:main_theorem_map}\label{tab:hypothesis_ledger}
\begin{tabular}{>{\raggedright\arraybackslash}p{0.30\linewidth}
                >{\raggedright\arraybackslash}p{0.33\linewidth}
                >{\raggedright\arraybackslash}p{0.26\linewidth}}
\toprule
Category and conclusion & Inputs & Status \\
\midrule
Finite-dimensional boundary theory: relative boundary Hamiltonian and collar comparison & Finite
triangulated surface with full
boundary subcomplex & Internal, \Cref{thm:relative_boundary_comparison} \\
Finite-dimensional jamming theory: finite-\(M\) decoupling, leakage, contour projections, Morse
index & Finite block separation
hypotheses & Internal, Section~\ref{sec:decoupling} \\
Internal bulk spectral measures: PL model, doubled compression, and finite-\(M\) jammed model have
the same normalized bulk spectral limits & Uniform finite-dimensional bounds, \(M\to\infty\), and
boundary-collar dimension \(o(\dim\Hil(\K_n))\) & Internal,
\Cref{thm:internal_bulk_spectral_equivalence} \\
Roe-algebra \(K_1\)-class identity: global interface
\(\partial_Y(E_+(\pi(P_\pm)))\in K_1(C^*(Y\subset X))\) & The Higson--Prodan closed-bulk
gap/Chern theorem and the internal original/copy Roe/Ewert--Meyer identification & Internal
modulo HP, \Cref{thm:pathb_roe_winding} \\
Toeplitz integer winding \(\pm1\) for straight periodic interfaces & Periodic triangular gap,
Chern signs, BDF--Toeplitz extension & Internal,
\Cref{cor:toeplitz_winding,thm:periodic_toeplitz_pairing} \\
Integer winding for general partitions & Hypothesized Chern-jump-compatible integer pairing on
\(K_1(\mathcal I)\) & Conditional on the pairing's existence,
\Cref{cor:conditional_winding} \\
Internal periodic examples: triangular closed-bulk input & Periodic triangular Bloch symbol &
Internal gap and Chern
signs, \Cref{thm:periodic_triangular_bulk,thm:periodic_triangular_chern} \\
\bottomrule
\end{tabular}
\end{table}

\begin{table}[H]
\centering
\small
\caption{Notation used for the bulk, jammed, and sidewise interface operators.}
\label{tab:notation_ledger}
\begin{tabular}{>{\raggedright\arraybackslash}p{0.24\linewidth}
                >{\raggedright\arraybackslash}p{0.65\linewidth}}
\toprule
Symbol & Meaning \\
\midrule
\(H_\pm=(B+B^\dagger)\pm S\) & Closed universal Chern Hamiltonian. \\
\(H_{\pm,M}=H_\pm(D(\K))+M\Proj_{\mathrm{copy}}\) & Finite doubled-and-jammed Hamiltonian. \\
\(A=P_{\mathrm{orig}}H_\pm(D(\K))P_{\mathrm{orig}}\) & Original-side compression of the doubled
Hamiltonian. \\
\(H_\pm^{\mathrm{PL}}(\K)\) & Intrinsic Poincar\'e--Lefschetz boundary Hamiltonian. \\
\(P_\pm=\chi_{(-\infty,0)}(H_\pm^\infty)\) & Closed-bulk Fermi projection in the coarse disjoint
union. \\
\(J_{\pm,M}^\infty\) & Direct-sum jammed family defining the interface. \\
\(T^+\), \(T^-\) & Positive-side and negative-side comparison operators. \\
\(Y\), \(C^*(Y\subset X)\) & Coarse interface and its Roe ideal. \\
\(\partial_Y(E_+(\pi(P_\pm)))\) & Internally constructed interface \(K_1\)-class. \\
\bottomrule
\end{tabular}
\end{table}

%====================================================================
\section{Simplicial setup and universal operators}\label{sec:setup}
%====================================================================

\begin{definition}[Triangulated surface with boundary]\label{def:surfwithbdry}
A finite oriented $2$-dimensional simplicial complex $\K$ is a \emph{triangulated surface with
boundary} if its geometric realization $|\K|$ is a compact PL $2$-manifold with boundary and the
boundary subcomplex $\bd\K$ is a \emph{full subcomplex}: every simplex of $\K$ whose vertices all
lie on $\bd\K$ is itself a boundary simplex. Every edge is incident to exactly one or two faces;
the edges incident to exactly one face form $\bd\K$, and $\bd\K$ is a disjoint union of cycles.
\end{definition}

The fullness condition is mild: it holds automatically after one barycentric subdivision. Its
role in keeping the simplicial double an abstract simplicial complex (without duplicate simplices
or boundary-edge valence failures) is discussed in detail in the supplement,
Section~\ref{S-sec:setup_remarks}.

We work on the simplicial Hilbert space
\(\Hil(\K)=\bigoplus_{p=0}^2 \ell^2(\K_p)\)
with the canonical orthonormal basis indexed by oriented simplices. The simplicial boundary
operator $B$ is the standard chain boundary, viewed as a degree-lowering operator. Then
$B^2=0$ and $B+B^\dagger$ is self-adjoint and local. We endow the set of barycenters
$X(\K)$ with the graph metric on the $1$-skeleton of the first barycentric subdivision; an
operator $T$ has \emph{propagation at most $R$} when
$\langle\sigma,T\tau\rangle=0$ for $\dist(\sigma,\tau)>R$.

For a \emph{closed} oriented triangulation $\K$, cap product with the fundamental cycle
$|X\rangle=\sum_{\tau\in\K_2}|\tau\rangle$ gives the Higson--Prodan operator $P$ from which one
defines a Hermitian, degree-reversing, local operator $S$ by the degree-dependent
symmetrization
\[
T|_{\Hil_p}=\tfrac12\bigl(P^\dagger+(-1)^pP\bigr),\qquad
S|_{\Hil_p}=i^{p(p-1)+1}T|_{\Hil_p},
\]
so that $BS+SB^\dagger=0$. The closed universal Chern Hamiltonian is
\(H_\pm(\K)=(B+B^\dagger)\pm S\).

\subsection{Closed-case gap criteria}
Two finite-dimensional gap criteria are used as building blocks; their proofs are short and
collected in the supplement.

\begin{proposition}[Exact gap lower bound under exact anticommutation]\label{prop:gap_from_S}
Let $\K$ be closed and assume $(B+B^\dagger)S+S(B+B^\dagger)=0$. Then
$H_\pm(\K)^2=(B+B^\dagger)^2+S^2$. If $\mu_S:=\dist(0,\spec(S))>0$, then
$\spec(H_\pm(\K))\cap(-\mu_S,\mu_S)=\emptyset$.
\end{proposition}
\proofsuppref{app:proof:prop:gap_from_S}

\begin{proposition}[Gap lower bound with anticommutation defect]\label{prop:gap_defect}
With $Q=B+B^\dagger$, $\mu_S:=\dist(0,\spec(S))$, and $\e_{\mathrm{ac}}:=\opnorm{QS+SQ}$, if
$S$ is Hermitian then $H_\pm(\K)^2\succeq(\mu_S^2-\e_{\mathrm{ac}})\Id$.
\end{proposition}
\proofsuppref{app:proof:prop:gap_defect}

\subsection{Doubling, jamming, and the boundary Hamiltonians}
Given $\K$ with boundary in the sense of Definition~\ref{def:surfwithbdry}, let $\overline\K$ be an
oriented copy with reversed face orientation, sharing each boundary vertex of $\K$. The
\emph{double} $D(\K)=\K\cup_{\bd\K}\overline\K$ is a finite closed oriented triangulated surface:

\begin{proposition}[The double is a closed oriented triangulated surface]\label{prop:double_closed}
Let $\K$ satisfy Definition~\ref{def:surfwithbdry}. Then $D(\K)$ is a finite closed oriented
abstract simplicial complex that triangulates a closed surface.
\end{proposition}
\proofsuppref{app:proof:prop:double_closed}

\begin{theorem}[Simplex count identities under doubling]\label{thm:counts}
\(|D(\K)_2|=2|\K_2|\),\, \(|D(\K)_1|=2|\K_1|-|\bd\K_1|\),\, \(|D(\K)_0|=2|\K_0|-|\bd\K_0|\), and
$|\bd\K_0|=|\bd\K_1|$. Consequently
$\dim\Hil(D(\K))=2\dim\Hil(\K)-2|\bd\K_0|$, $\chi(D(\K))=2\chi(\K)$, and for a connected,
orientable, genus-$g$ surface with $b\ge 1$ boundary components, the doubled genus is
$g_{\mathrm{double}}=2g+b-1$.
\end{theorem}
\proofsuppref{app:proof:thm:counts}

The closed-case construction therefore produces the operator $H_\pm(D(\K))$ on $\Hil(D(\K))$.
We split
\[
\Hil(D(\K))=\Hil_{\mathrm{orig}}\oplus\Hil_{\mathrm{copy}}
\]
and \emph{jam} the reflected copy by the on-site potential
\begin{equation}\label{eq:Hjam}
H_{\pm,M}=H_\pm(D(\K))+M\Proj_{\mathrm{copy}}.
\end{equation}
The block decomposition
\(H_\pm(D(\K))=\begin{pmatrix}A&C\\ C^\dagger&D\end{pmatrix}\)
on $\Hil_{\mathrm{orig}}\oplus\Hil_{\mathrm{copy}}$ defines the
\emph{original-side boundary Hamiltonian}
\(H_{\pm}^{\bd}(\K):=A\), a self-adjoint local operator on the surface-with-boundary.

The intrinsic alternative uses the relative chain complex
$C_p(\K,\bd\K)=C_p(\K)/C_p(\bd\K)$ and the relative fundamental chain
$[\K,\bd\K]=\sum_{\tau\in\K_2}|\tau\rangle$. Cap product with $[\K,\bd\K]$, followed by the
canonical lift $\widehat P_p:=j_{2-p}P^{\mathrm{rel}}_p:C_p(\K)\to C_{2-p}(\K)$, gives the
Hermitian, degree-reversing, local relative duality operator
$S^{\mathrm{PL}}_\K$ on $\Hil(\K)$. The \emph{relative Poincar\'e--Lefschetz boundary
Hamiltonian} is
\begin{equation}\label{eq:H_PL}
H^{\mathrm{PL}}_{\pm}(\K):=B_\K+B_\K^\dagger\pm S^{\mathrm{PL}}_\K\quad\text{on }\Hil(\K).
\end{equation}

\begin{lemma}[Relative cap identity and collar defect]\label{lem:relative_cap_identity}
With $\partial^{\mathrm{rel}}_r:C_r(\K,\bd\K)\to C_{r-1}(\K,\bd\K)$ defined by
$\partial^{\mathrm{rel}}_r q_r=q_{r-1}B_{\K,r}$, and $\delta_p:=B_{\K,p+1}^\dagger$, the
relative cap-product blocks satisfy
$\partial^{\mathrm{rel}}_{2-p}P^{\mathrm{rel}}_p+(-1)^p P^{\mathrm{rel}}_{p+1}\delta_p=0$
for $p=0,1$. After lifting by $j$, the difference
\(B_{\K,2-p}\widehat P_p+(-1)^p\widehat P_{p+1}\delta_p\) has matrix coefficients only when
both source and target lie in a fixed barycentric collar of $\bd\K$. Consequently the relative
duality defect $B_\K S^{\mathrm{PL}}_\K+S^{\mathrm{PL}}_\K B_\K^\dagger$ is supported in a fixed
boundary collar.
\end{lemma}
\proofsuppref{app:proof:lem:relative_cap_identity}

\begin{lemma}[Coefficient agreement off the boundary collar]\label{lem:cap_coefficient_agreement}
Let $P_D$ be the closed cap-product operator on $D(\K)$ and $\widehat P$ the lifted relative
cap-product operator on $\K$. There is a radius $R_\frown$ depending only on the local
cap-product rule, such that
$\langle\sigma,P_D\tau\rangle=\langle\sigma,\widehat P\tau\rangle$ whenever
$\sigma,\tau\in\K\setminus N_{R_\frown}(\bd\K)$.
\end{lemma}
\proofsuppref{app:proof:lem:cap_coefficient_agreement}

\begin{theorem}[Relative boundary Hamiltonian and doubled compression]\label{thm:relative_boundary_comparison}
Let \(\K\) be a triangulated oriented surface with full boundary subcomplex. Under the natural
identification of \(\Hil(\K)\) with \(\Hil_{\mathrm{orig}}\subset\Hil(D(\K))\),
\(H^{\mathrm{PL}}_{\pm}(\K)\) is self-adjoint and local. Moreover there is a self-adjoint operator
\(R_{\partial,\pm}\) such that
\begin{equation}\label{eq:relative_boundary_comparison}
H_{\pm}^{\bd}(\K)
=
H^{\mathrm{PL}}_{\pm}(\K)+R_{\partial,\pm},
\end{equation}
and the comparison error \(R_{\partial,\pm}\), and also the relative duality defect
\(B_\K S^{\mathrm{PL}}_\K+S^{\mathrm{PL}}_\K B_\K^\dagger\), are supported in a bounded collar of
the boundary: there is a radius \(R_{\mathrm{PL}}\), depending only on the local cap-product
convention, such that the matrix entries vanish unless both source and target lie in
\(N_{R_{\mathrm{PL}}}(\bd\K)\). For uniformly regular refinement families the same
\(R_{\mathrm{PL}}\) works for all members.
\end{theorem}
\proofsuppref{app:proof:thm:relative_boundary_comparison}

\begin{lemma}[Spectral counts under collar-supported perturbations]\label{lem:collar_rank_counts}
If $R=P_W R P_W$ for a projection $P_W$ of rank $r$, then
$|N_{A+R}(t)-N_A(t)|\le r$ for all $t$ and
$|\rank\chi_I(A+R)-\rank\chi_I(A)|\le 2r$ for every interval $I$. Applied to
$A=H^{\mathrm{PL}}_{\pm}(\K)$ and $A+R=H_{\pm}^{\bd}(\K)$, one may take
$W=\ell^2(N_{R_{\mathrm{PL}}}(\bd\K))$. In particular, along refinement families with
$r_n/d_n\to0$, the normalized spectral distribution functions of the two boundary models share
the same subsequential limits.
\end{lemma}
\proofsuppref{app:proof:lem:collar_rank_counts}

%====================================================================
\section{Finite-dimensional decoupling theory}\label{sec:decoupling}
%====================================================================

The finite-dimensional analysis is perturbation theory specialized to the block structure produced
by doubling and jamming. We use standard Schur-complement and spectral projection tools
\cite[Chapter~II]{KatoPerturbation,BhatiaMatrixAnalysis}, but keep the constants and separation
hypotheses explicit because they are later checked on the computed matrices. The full proofs are
collected in the online supplement, Section~\ref{S-sec:schur_proofs}.

\begin{lemma}[Schur complement resolvent formula]\label{lem:schur}
Let $T=\bigl(\begin{smallmatrix}A&C\\ C^\dagger&D\end{smallmatrix}\bigr)$ be self-adjoint on
$\Hil_1\oplus\Hil_2$ and $z\in\mathbb C\setminus\mathbb R$. Then the $(1,1)$ block of $(T-z)^{-1}$
is $\bigl(A-z-C(D-z)^{-1}C^\dagger\bigr)^{-1}$.
\end{lemma}

Write $D_M:=D+M\Id$ and $d_M(z):=\dist(z,\spec(D_M))$.

\begin{theorem}[Compressed resolvent convergence with explicit $O(1/M)$ rate]\label{thm:resolvent}
Let $H_{\pm,M}$ be as in \eqref{eq:Hjam}. For any $z\in\mathbb C\setminus\mathbb R$,
\begin{equation}\label{eq:resolvent_bound_exact}
\opnorm{P_{\mathrm{orig}}(H_{\pm,M}-z)^{-1}P_{\mathrm{orig}}-(A-z)^{-1}}
\le
\frac{\opnorm{C}^2}{d_M(z)\,|\Im z|^2}.
\end{equation}
In particular, if $M>\opnorm{D}+|\Re z|$,
$P_{\mathrm{orig}}(H_{\pm,M}-z)^{-1}P_{\mathrm{orig}}\to(A-z)^{-1}$ in operator norm at rate
$O(1/M)$, so the doubled-and-jammed Hamiltonian recovers the original-side boundary Hamiltonian
$A=H_{\pm}^{\bd}(\K)$ in the norm-resolvent sense as $M\to\infty$.
\end{theorem}

\begin{proof}[Proof sketch; full proof in Supplement, Section~\ref{S-sec:schur_proofs}.]
Write $R_M:=P_{\mathrm{orig}}(H_{\pm,M}-z)^{-1}P_{\mathrm{orig}}$ and $R_0:=(A-z)^{-1}$.
The Schur complement (Lemma~\ref{lem:schur}) gives the resolvent identity
$R_M-R_0=R_M\,K_M\,R_0$ with the self-energy correction
$K_M:=C(D_M-z)^{-1}C^\dagger$. The resolvent norms satisfy
$\opnorm{R_M}\le|\Im z|^{-1}$ and $\opnorm{R_0}\le|\Im z|^{-1}$, and the self-energy is
bounded by $\opnorm{K_M}\le\opnorm{C}^2/d_M(z)$. Multiplying gives
\eqref{eq:resolvent_bound_exact}.
\end{proof}

\begin{corollary}[Boundary Hamiltonian as the large-jam limit]\label{cor:boundary_hamiltonian_limit}
Let $H_{\pm}^{\bd}(\K)=A$ be the original-side boundary Hamiltonian. For every
$z\in\mathbb C\setminus\mathbb R$,
\(P_{\mathrm{orig}}(H_{\pm,M}-z)^{-1}P_{\mathrm{orig}}\to(H_{\pm}^{\bd}(\K)-z)^{-1}\) in operator
norm with the explicit $O(M^{-1})$ bound \eqref{eq:resolvent_bound_exact}. Thus the
doubled-and-jammed Hamiltonian recovers a self-adjoint boundary Hamiltonian on the original side
in the norm-resolvent sense; equivalently, by Theorem~\ref{thm:relative_boundary_comparison},
the same limit is $(H^{\mathrm{PL}}_{\pm}(\K)+R_{\partial,\pm}-z)^{-1}$.
\end{corollary}
\proofsuppref{app:proof:cor:boundary_hamiltonian_limit}

\begin{corollary}[Degenerate full block-resolvent limit]\label{cor:full_resolvent}
For $z\in\mathbb C\setminus\mathbb R$, $(H_{\pm,M}-z)^{-1}\to(A-z)^{-1}\oplus0$ in operator norm
at rate $O(1/M)$, with the explicit triangle inequality
\(\|(H_{\pm,M}-z)^{-1}-((A-z)^{-1}\oplus0)\|\le \opnorm{C}^2/(d_M(z)|\Im z|^2)
+ 2\opnorm{C}/(d_M(z)|\Im z|)+1/d_M(z)+\opnorm{C}^2/(d_M(z)^2|\Im z|)\).
This is generalized norm-resolvent convergence in which the copy sector is sent to infinite energy.
\end{corollary}
\proofsuppref{app:proof:cor:full_resolvent}

\begin{theorem}[First-order Schur-complement expansion]\label{thm:first_order_eff}
Fix $z\in\mathbb C\setminus\mathbb R$. If $M>\opnorm{D-z}$ and $A_M^{(1)}:=A-\frac1MCC^\dagger$,
then
\begin{equation}\label{eq:first_order_bound}
\opnorm{P_{\mathrm{orig}}(H_{\pm,M}-z)^{-1}P_{\mathrm{orig}}-(A_M^{(1)}-z)^{-1}}
\le
\frac{\opnorm{C}^2\,\opnorm{D-z}}{M(M-\opnorm{D-z})\,|\Im z|^2}=O(M^{-2}).
\end{equation}
\end{theorem}
\proofsuppref{app:proof:thm:first_order_eff}

For a positively oriented simple closed contour $\Gamma\subset\rho(T)$ write the Riesz projection
$P_\Gamma(T):=\frac{1}{2\pi i}\int_\Gamma (z-T)^{-1}\,dz$. Set
$\delta_A(\Gamma):=\dist(\Gamma,\spec(A))$,
$\delta_D(\Gamma):=\dist(\Gamma,\spec(D_M))$, and
$\delta_H(\Gamma):=\delta_A(\Gamma)-\opnorm{C}^2/\delta_D(\Gamma)$.

\begin{theorem}[General contour-projection comparison]\label{thm:proj_general}
Assume $\delta_A(\Gamma)>0$, $\delta_D(\Gamma)>0$, $\delta_H(\Gamma)>0$. Then
$\Gamma\subset\rho(H_{\pm,M})$,
$\rank P_\Gamma(H_{\pm,M})=\rank P_\Gamma(A)+\rank P_\Gamma(D_M)$, and
\begin{equation}\label{eq:proj_general_bound}
\opnorm{P_\Gamma(H_{\pm,M})-\bigl(P_\Gamma(A)\oplus P_\Gamma(D_M)\bigr)}
\le \tfrac{|\Gamma|}{2\pi}\,\eta_M(\Gamma)=O(M^{-1}),
\end{equation}
where $\eta_M(\Gamma)$ is the explicit constant
\[
\eta_M(\Gamma)
=\frac{\opnorm C^2}{\delta_A(\Gamma)\delta_H(\Gamma)\delta_D(\Gamma)}
+\frac{2\opnorm C}{\delta_H(\Gamma)\delta_D(\Gamma)}
+\frac{\opnorm C^2}{\delta_H(\Gamma)\delta_D(\Gamma)^2}.
\]
\end{theorem}
\proofsuppref{app:proof:thm:proj_general}

For $E_{\max}>0$ define $g_M(E_{\max}):=\dist\bigl([-E_{\max},E_{\max}],\spec(D_M)\bigr)$.

\begin{theorem}[Bounded-energy spectral localization]\label{thm:spectral_localization}
If $E$ is an eigenvalue of $H_{\pm,M}$ with $|E|\le E_{\max}$ and $g_M(E_{\max})>0$, then
$\dist(E,\spec(A))\le\opnorm{C}^2/g_M(E_{\max})$.
\end{theorem}
\proofsuppref{app:proof:thm:spectral_localization}

\begin{corollary}[Spectral comparison with the relative boundary Hamiltonian]\label{cor:relative_spectral_localization}
Under the hypotheses of Theorem~\ref{thm:spectral_localization},
\(\dist\bigl(E,\spec(H^{\mathrm{PL}}_\pm(\K)+R_{\partial,\pm})\bigr)
\le\opnorm{C}^2/g_M(E_{\max})\).
\end{corollary}
\proofsuppref{app:proof:cor:relative_spectral_localization}

\begin{theorem}[Bounded-energy eigenvalues lie near the first-order effective operator]\label{thm:first_order_spectral_localization}
Let $E$ be an eigenvalue of $H_{\pm,M}$ with $|E|\le E_{\max}$ and $g_M(E_{\max})>0$. Then
\begin{equation}\label{eq:specA1_exact}
\dist\bigl(E,\spec(A_M^{(1)})\bigr)
\le
\frac{\opnorm{C}^2\,(\opnorm{D}+E_{\max})}{M\,g_M(E_{\max})}.
\end{equation}
In particular, if $M>\opnorm{D}+E_{\max}$, then
\begin{equation}\label{eq:specA1_crude}
\dist\bigl(E,\spec(A_M^{(1)})\bigr)
\le
\frac{\opnorm{C}^2\,(\opnorm{D}+E_{\max})}{M(M-\opnorm{D}-E_{\max})}.
\end{equation}
\end{theorem}
\proofsuppref{app:proof:thm:first_order_spectral_localization}

\begin{theorem}[Copy-side leakage bound]\label{thm:eigvec_leak}
If $(E,\psi)$ is an eigenpair of $H_{\pm,M}$ with $|E|\le E_{\max}$ and $g_M(E_{\max})>0$,
the copy component satisfies $\|\psi_{\mathrm{copy}}\|\le
\opnorm{C}\,\|\psi_{\mathrm{orig}}\|/g_M(E_{\max})=O(1/M)$.
\end{theorem}
\proofsuppref{app:proof:thm:eigvec_leak}

\begin{theorem}[Inertia reduction for large scalar jamming]\label{thm:inertia_reduction}
Assume $M>\opnorm{D}$, so $D_M=D+M\Id$ is positive, and set $S_M(0):=A-CD_M^{-1}C^\dagger$. Then
$N_-(H_{\pm,M})=N_-(S_M(0))$, $\dim\ker H_{\pm,M}=\dim\ker S_M(0)$, and
$\|S_M(0)-A\|\le\opnorm{C}^2/(M-\opnorm{D})$. Consequently, if $0\notin\spec(A)$ and the
perturbation satisfies $\opnorm{C}^2/(M-\opnorm{D})<\dist(0,\spec(A))$, then $H_{\pm,M}$ is
invertible and $N_-(H_{\pm,M})=N_-(A)$.
\end{theorem}
\proofsuppref{app:proof:thm:inertia_reduction}

\begin{corollary}[Monotonicity in the scalar jam scale]\label{cor:scalar_jam_monotonicity}
Let $M_2>M_1$. The eigenvalues $\lambda_j(H_{\pm,M})$, listed in nondecreasing order with
multiplicity, are nondecreasing in $M$. In particular, $M\mapsto N_-(H_{\pm,M})$ is nonincreasing.
\end{corollary}
\proofsuppref{app:proof:cor:scalar_jam_monotonicity}

\begin{corollary}[No low-energy pure-copy states]\label{cor:no_copy_spectral_subspace}
Let $E_{\max}>0$ and assume $M>\opnorm{D}+E_{\max}$. The compression
$P_{\mathrm{orig}}\,Q_M:\Ran Q_M\to\Hil_{\mathrm{orig}}$ is injective for
$Q_M=\chi_{[-E_{\max},E_{\max}]}(H_{\pm,M})$, so $\rank Q_M\le\dim\Hil_{\mathrm{orig}}$, and
quantitatively $\|P_{\mathrm{copy}}\psi\|\le (E_{\max}+\opnorm C)/(M-\opnorm D)\,\|\psi\|$ for
every $\psi\in\Ran Q_M$.
\end{corollary}
\proofsuppref{app:proof:cor:no_copy_spectral_subspace}

\begin{proposition}[Inherited gap from the original-side block]\label{prop:inherited_gap_A}
Assume $\spec(A)\cap(-g_A,g_A)=\emptyset$ for some $g_A>0$ and
$g_M(g_A):=\dist([-g_A,g_A],\spec(D_M))>0$. With $\Delta_M=\opnorm{C}^2/g_M(g_A)$, if
$\Delta_M<g_A$ then $\spec(H_{\pm,M})\cap(-g_A+\Delta_M,g_A-\Delta_M)=\emptyset$.
\end{proposition}
\proofsuppref{app:proof:prop:inherited_gap_A}

\begin{corollary}[Uniform direct-sum large-jam estimates]\label{cor:uniform_direct_sum_jam}
Under uniform bounds $\bar C=\sup_n\opnorm{C_n}<\infty$ and $\bar D=\sup_n\opnorm{D_n}<\infty$,
the finite-$M$ estimates of
\Cref{thm:resolvent,thm:first_order_eff,thm:spectral_localization,thm:eigvec_leak} hold uniformly
in $n$ for the direct-sum doubled-and-jammed family. The compressed resolvent error, the
first-order compressed approximation, the spectral-localization distance, and the copy leakage
carry the same explicit bounds with $\opnorm{C}$ replaced by $\bar C$ and $\opnorm{D}$ by
$\bar D$.
\end{corollary}
\proofsuppref{app:proof:cor:uniform_direct_sum_jam}

\begin{theorem}[Internal bulk spectral equivalence]\label{thm:internal_bulk_spectral_equivalence}
Let $(\K_n)_n$ be finite triangulated oriented surfaces with full boundary subcomplex, with
$d_n=\dim\Hil(\K_n)$ and uniform bounds on $\|A_{\pm,n}\|,\|H^{\mathrm{PL}}_\pm(\K_n)\|,
\|C_{\pm,n}\|,\|D_{\pm,n}\|$. Let $M_n\to\infty$ and assume $r_n/d_n\to0$ for the
boundary-collar dimension $r_n$ from Theorem~\ref{thm:relative_boundary_comparison}. Then for
every $f\in C_0(\mathbb R)$,
\[
\frac1{d_n}\operatorname{Tr} f(H_{\pm,n,M_n})
-
\frac1{d_n}\operatorname{Tr} f\!\left(H^{\mathrm{PL}}_{\pm}(\K_n)\right)
\longrightarrow0.
\]
The intrinsic Poincar\'e--Lefschetz model, the doubled original-side compression, and the
finite-\(M_n\) jammed model have the same normalized bulk spectral subsequential limits. The
statement uses no closed-bulk Chern theorem, no uniform doubled-bulk gap, and no
interface-pairing hypothesis.
\end{theorem}
\proofsuppref{app:proof:thm:internal_bulk_spectral_equivalence}

The remaining estimates in this section -- contour eigenvalue-count stability, monotonicity in the
jam scale, no-copy low-energy bounds, an inherited-gap criterion, and the uniform direct-sum
version used by the Roe construction -- are stated together with their explicit constants and
proved in the supplement, Section~\ref{S-sec:schur_proofs}.

%====================================================================
\section{Roe-algebra identification of the doubled-and-jammed interface class}\label{sec:pathb_internalization}
%====================================================================

This section establishes the boundary-map identification used by the headline Roe-winding theorem.
The closed-bulk gap and Chern value come from the closed universal-model theorem of Higson and
Prodan \cite{HigsonProdanUniversalPRL,HigsonProdanUniversalArxiv}. The Ewert--Meyer citation
\cite{EwertMeyer2019} fixes the coarse bulk--interface language and the Mayer--Vietoris pairing
convention; for the doubled-and-jammed original/copy partition used in this paper, the
identification of the Mayer--Vietoris boundary with the Roe connecting map is proved here.

\subsection{Higson--Prodan universal closed-bulk input}

\begin{definition}[Universal Hamiltonian on a bounded-geometry triangulation]\label{def:pathb_universal_hamiltonian}
Let \(K\) be a bounded-geometry closed oriented triangulated surface, let
\(\Hil(K)=\ell^2(K_0\sqcup K_1\sqcup K_2)\), let \(B\) be the simplicial boundary operator, and let
\(S_K\) be the Higson--Prodan Hermitian Poincar\'e-duality operator constructed from cap product
with the fundamental cycle. The two universal closed Hamiltonians are
\[
H_\pm(K):=(B+B^\dagger)\pm S_K
\]
on \(\Hil(K)\).
\end{definition}

\begin{lemma}[Higson--Prodan closed-bulk gap, cited]\label{lem:pathb_hp_gap}
For the universal Hamiltonians \(H_\pm(K)\) on closed bounded-geometry triangulations covered by
the Higson--Prodan universal-model theorem, there is a uniform Fermi gap
\(\spec(H_\pm(K))\cap(-g_{\mathrm{HP}},g_{\mathrm{HP}})=\varnothing\) for some
\(g_{\mathrm{HP}}>0\) independent of the refinement level.
\end{lemma}

\begin{lemma}[Higson--Prodan closed-bulk Chern number, cited]\label{lem:pathb_hp_chern}
With \(P_\pm(K)=\chi_{(-\infty,0)}(H_\pm(K))\), the Bellissard--Connes/Higson--Prodan closed-bulk
Chern pairing satisfies \(\Ch([P_\pm])=\pm1\), with the sign determined by the chosen component
and by the orientation convention of Section~\ref{sec:setup}.
\end{lemma}

\subsection{Roe extension and the interface boundary class}

Let \(X=\bigsqcup_n X(D(\K_n))\) be the coarse disjoint union of the doubled barycentric vertex
sets, let \(X^+\) be the original side, let \(X^-\) be the reflected-copy side, and let \(Y\) be a
fixed coarse collar of the glued boundary. The choice of coarse collar radius does not affect the
ideal $C^*(Y\subset X)$ (Proposition~\ref{S-prop:collar_radius_independence}).

\begin{lemma}[Roe algebra and closed-bulk Fermi projection]\label{lem:pathb_roe_setup}
For \(X\) as above, the scalar Roe algebra \(C^*(X)\) is unital. Under the Higson--Prodan gap of
Lemma~\ref{lem:pathb_hp_gap}, the direct-sum Fermi projection
\(P_\pm=\chi_{(-\infty,0)}\!\left(\bigoplus_n H_\pm(D(\K_n))\right)\) belongs to \(C^*(X)\).
\end{lemma}

\begin{lemma}[Coarse partition exact sequence]\label{lem:pathb_exact_sequence}
With \(\mathcal A=C^*(X)\) and \(\mathcal I=C^*(Y\subset X)\), the original/copy decomposition gives
the quotient extension
\(0\to\mathcal I\to\mathcal A\xrightarrow{\pi}\mathcal A/\mathcal I\to0\), with boundary map
\(\partial_Y:K_0(\mathcal A/\mathcal I)\to K_1(\mathcal I)\).
\end{lemma}

\begin{definition}[Interface projection and unitary]\label{def:pathb_interface_projection}
Let \(J_{\pm,M}^\infty=\bigoplus_n H_{\pm,M,n}\) be the doubled-and-jammed family and let
\(p_J:=E_+\bigl(\pi(P_\pm)\bigr)\in\mathcal A/\mathcal I\). For any admissible cutoff \(h\),
\(p_J=\pi(h(J_{\pm,M}^\infty))\), and the associated interface \(K_1\)-class is represented by
\(u_{J,h}:=\exp(2\pi i\,h(J_{\pm,M}^\infty))\in\Id+\mathcal I\),
\([u_{J,h}]=\partial_Y([p_J])\).
\end{definition}

\subsection{Conventions for Theorem~\ref{thm:pathb_roe_winding}}\label{sec:pathb_conventions}
The headline theorem fixes one set of conventions, used identically in the body proof and the
appendix-level verifications:
\begin{enumerate}[label=(C\arabic*),leftmargin=2.4em]
\item\label{conv:pathb_rep} \emph{Representation.} The standard left-regular
\(\ell^2\)-representation of the scalar Roe algebra: \(C^*(X)\subset\mathcal B(\ell^2(X))\) is the
norm closure of locally compact finite-propagation operators acting on the canonical basis of unit
point masses at each vertex of $X=\bigsqcup_n X(D(\K_n))$.
\item\label{conv:pathb_stab} \emph{Stabilization.} The scalar (not matrix-stabilized) Roe algebra.
Standard Murray--von Neumann \(K\)-theoretic stabilization is invoked only where strictly
necessary, namely in step (3) of the proof below.
\item\label{conv:pathb_sign} \emph{Sign convention for \(\partial\).} The connecting map
\(\partial_Y\) is the standard Higson--Roe boundary map in the six-term exact sequence of
\(0\to\mathcal I\to\mathcal A\to\mathcal A/\mathcal I\to0\), with the orientation choice that
\(\partial[\pi(P)]=[\exp(2\pi i\,\tilde P)]\) for any self-adjoint lift \(\tilde P\) of the
quotient projection \cite{HigsonRoeI,RoeLectures}.
\item\label{conv:pathb_side} \emph{Side convention.} \(X=X^+\sqcup X^-\) is the original/copy
decomposition; spectral cutoffs are taken on the negative real axis,
\(P_\pm=\chi_{(-\infty,0)}(H_\pm^\infty)\).
\item\label{conv:pathb_quotient} \emph{Quotient map.} \(\pi:\mathcal A\to\mathcal A/\mathcal I\) is
the canonical \(C^*\)-quotient by \(\mathcal I=C^*(Y\subset X)\), with collar-radius independence.
\end{enumerate}
The Ewert--Meyer Mayer--Vietoris boundary map for the same original/copy partition is denoted
\(\partial^{\mathrm{EM}}_Y\); the internal compatibility theorem below identifies it with
\(\partial_Y\) under \ref{conv:pathb_rep}--\ref{conv:pathb_quotient}.

\begin{figure}[H]
\centering
\[
\begin{tikzcd}[row sep=2.6em, column sep=2.6em]
0 \arrow[r] & \mathcal I \arrow[r, hook] \arrow[d, equal] & \mathcal A
  \arrow[r, two heads, "\pi"] \arrow[d, equal] & \mathcal A/\mathcal I \arrow[r]
  \arrow[d, equal] & 0 \\
0 \arrow[r] & C^*(Y\subset X) \arrow[r, hook] & C^*(X)
  \arrow[r, two heads, "\pi"] & C^*(X)/C^*(Y\subset X) \arrow[r] & 0 \\
K_1(\mathcal I) \arrow[r, "\cong"] & {[u_{J,h}]}
  & {[P_\pm]\in K_0(\mathcal A)} \arrow[r, "\pi_*"] \arrow[d, "\Ch"', "\mathrm{HP}"]
  & {[p_J]=[E_+(\pi(P_\pm))]\in K_0(\mathcal A/\mathcal I)}
    \arrow[ll, bend right=18, dashed, "\partial_Y\,=\,\partial^{\mathrm{EM}}_Y"'] \\
  & & \mathbb Z\ni\pm1 &
\end{tikzcd}
\]
\caption{Commutative diagram for Theorem~\ref{thm:pathb_roe_winding}. The dashed arrow labelled
\(\partial_Y=\partial^{\mathrm{EM}}_Y\) records the internal compatibility theorem: the Roe
connecting map and the Ewert--Meyer Mayer--Vietoris boundary coincide for the original/copy
partition, sending \([p_J]\) to the interface \(K_1\)-class
\([u_{J,h}]=[\exp(2\pi i\,h(J_{\pm,M}^\infty))]\). The downward arrow \(\mathrm{HP}\) is
Higson--Prodan's closed-bulk Chern character, which evaluates the bulk class to \(\pm1\).}
\label{fig:pathb_diagram}
\end{figure}

\subsection{Internal compatibility theorem}

\begin{theorem}[Roe interface class from the Higson--Prodan closed bulk]\label{thm:pathb_roe_winding}
Let \((\K_n)_n\) be a uniformly regular boundary refinement family whose closed doubles
\(D(\K_n)\) satisfy the Higson--Prodan universal-model theorem
\cite{HigsonProdanUniversalPRL,HigsonProdanUniversalArxiv}. Use the conventions
\ref{conv:pathb_rep}--\ref{conv:pathb_quotient}. Let \(\mathcal A=C^*(X)\),
\(\mathcal I=C^*(Y\subset X)\), and
\[
P_\pm=\chi_{(-\infty,0)}\!\left(\bigoplus_n H_\pm(D(\K_n))\right).
\]
For every \(M>B_{\mathrm{copy},-}\), the doubled-and-jammed family \(J_{\pm,M}^\infty\)
determines the quotient projection \(p_J=E_+(\pi(P_\pm))\in\mathcal A/\mathcal I\) and the
\(K_1\)-class identity
\[
\Indint(J_{\pm,M}^\infty;T^+,T^-)
=
\partial_Y(E_+(\pi(P_\pm)))
=
[\,\exp(2\pi i\,h(J_{\pm,M}^\infty))\,]\in K_1(\mathcal I)
\]
for any admissible cutoff \(h\). The Roe connecting map \(\partial_Y\) coincides on
\([p_J]\) with the Ewert--Meyer Mayer--Vietoris boundary map \(\partial^{\mathrm{EM}}_Y\)
attached to the original/copy partition under
\ref{conv:pathb_rep}--\ref{conv:pathb_quotient}; this identification is proved internally in
Step (2) of the body proof and is the dashed arrow of Figure~\ref{fig:pathb_diagram}. The
class \(\partial_Y(E_+(\pi(P_\pm)))\) is independent of the admissible cutoff \(h\) and of
the jam scale \(M>B_{\mathrm{copy},-}\), and is preserved under refinement subsequences for
which the Higson--Prodan Fermi gap is uniform.
\end{theorem}

\begin{proof}[Body proof of Theorem~\ref{thm:pathb_roe_winding}]
\textbf{Roadmap.} The proof has four steps, in the order displayed by
Figure~\ref{fig:pathb_diagram}. Steps (1), (2), and (4) are internal to this paper. Step (3) is
the only step that uses an external theorem, namely Higson--Prodan
\cite{HigsonProdanUniversalPRL,HigsonProdanUniversalArxiv}. The convention block
\ref{conv:pathb_rep}--\ref{conv:pathb_quotient} fixes the representation, stabilization, sign,
side, and quotient choices used uniformly across all four steps.

\medskip\noindent\textbf{Step (1): Roe boundary map equals the cutoff unitary.}
By Lemma~\ref{lem:pathb_roe_setup}, \(P_\pm\in\mathcal A=C^*(X)\) under the Higson--Prodan gap of
Lemma~\ref{lem:pathb_hp_gap}, so the quotient projection
\(p_J:=E_+(\pi(P_\pm))\in\mathcal A/\mathcal I\) is well defined. The functional calculus of the
quotient-sidewise statement (Proposition~\ref{S-prop:quotient_sidewise} in the supplement)
provides, for any admissible cutoff \(h\), a self-adjoint element
\(h(J_{\pm,M}^\infty)\in\mathcal A\) with \(\pi(h(J_{\pm,M}^\infty))=p_J\); see
Definition~\ref{def:pathb_interface_projection} for the explicit lift. Since
\(h(J_{\pm,M}^\infty)\) is a self-adjoint lift of an idempotent in the quotient,
\((h(J_{\pm,M}^\infty))^2-h(J_{\pm,M}^\infty)\in\mathcal I\), hence
\(\exp(2\pi i\,h(J_{\pm,M}^\infty))\in\Id+\mathcal I\). Convention~\ref{conv:pathb_sign} fixes the
sign of the standard Higson--Roe \(K\)-theoretic connecting map, and the boundary-map lemma
(Lemma~\ref{S-lem:pathb_boundary_map}) gives
\[
\partial_Y([p_J])
=
[\exp(2\pi i\,h(J_{\pm,M}^\infty))]\in K_1(\mathcal I),
\]
independent of the admissible cutoff. This is the left dashed arrow in
Figure~\ref{fig:pathb_diagram}.

\medskip\noindent\textbf{Step (2): Roe and Ewert--Meyer boundary maps coincide for the
original/copy partition.}
The coarse-excision statement (Proposition~\ref{S-prop:coarse_excision}) shows that the
original/copy decomposition \((X^+,X^-;Y)\) is coarsely excisive: the interface ideal
\(\mathcal I=C^*(Y\subset X)\) is the kernel of the canonical restriction map to the disjoint
quotient. Consequently the same six-term exact sequence underlies both the Roe-extension
connecting map \(\partial_Y\) and the Ewert--Meyer Mayer--Vietoris boundary map
\(\partial^{\mathrm{EM}}_Y\) associated to the bulk--interface pair \cite{EwertMeyer2019}. By
naturality of the connecting map, any two extensions identifying the same short-exact sequence of
\(C^*\)-algebras with the same orientation (fixed in Convention~\ref{conv:pathb_sign}) produce
the same \(K_1\)-class. The quotient-splitting identity
\(\mathcal A/\mathcal I=E_+(\mathcal A/\mathcal I)\oplus E_-(\mathcal A/\mathcal I)\)
(Proposition~\ref{S-prop:quotient_split}) identifies the positive and negative comparison
projections, with the negative side trivial by the trivial-jam statement
(Proposition~\ref{S-prop:trivial_jam}). The \(K_1\)-class identity in the theorem statement
follows from naturality of the connecting map applied to the coarsely excisive short-exact
sequence.

\medskip\noindent\textbf{Step (3): Closed-bulk Chern from Higson--Prodan.}
By Lemma~\ref{lem:pathb_hp_chern}, which paraphrases the Higson--Prodan universal-model theorem
\cite{HigsonProdanUniversalPRL,HigsonProdanUniversalArxiv} for closed bounded-geometry
triangulated surfaces, the closed-bulk Fermi projection has Chern number
\(\Ch([P_\pm])=\pm1\). This is the only step that uses an external theorem. The integer
evaluation of the \(K_1\)-class identity from Step (2) is taken up by the two corollaries
(\Cref{cor:toeplitz_winding} for the internal Toeplitz periodic case;
\Cref{cor:conditional_winding} for the general conditional case).

\medskip\noindent\textbf{Step (4): Finite-\(M\) quotient identification and \(K_1\)-limit.}
For every \(M>B_{\mathrm{copy},-}\), the trivial-jam statement
(Proposition~\ref{S-prop:trivial_jam}) gives strict positivity of the negative comparison operator,
so its negative Fermi projection vanishes. The positive comparison operator is
\(T^+=H_\pm^\infty\) with negative Fermi projection \(P_\pm\), and the quotient-sidewise
identification (Proposition~\ref{S-prop:quotient_sidewise}) gives
\(p_J=E_+(\pi(P_\pm))+E_-(\pi(0))=E_+(\pi(P_\pm))\) at every finite jam scale
\(M>B_{\mathrm{copy},-}\). The corresponding \(K_1\)-class is independent of \(M\)
(Corollary~\ref{S-cor:jam_independence}), so the limit \(M\to\infty\) is well defined as a
single \(K_1(\mathcal I)\)-class. Composing the four steps gives the claimed
\(K_1\)-class identity
\(\Indint(J_{\pm,M}^\infty;T^+,T^-)=\partial_Y(E_+(\pi(P_\pm)))
=[\exp(2\pi i\,h(J_{\pm,M}^\infty))]\). Detailed computational verifications of the
cutoff-functional-calculus identity, the norm-continuous homotopy along admissible cutoffs,
and the explicit lift \(h(J_{\pm,M}^\infty)\) are collected in the supplement,
Section~\ref{S-app:proof:thm:pathb_roe_winding}.
\end{proof}

\begin{corollary}[Integer winding in the Toeplitz periodic case]\label{cor:toeplitz_winding}
Suppose \(Y\subset X\) is the straight periodic interface of the periodic triangular class of
Section~\ref{sec:periodicbulk}, with primitive tangential/positive-normal basis \((t,n)\) and
\(\varepsilon(t,n)=\det[t\ n]\in\{\pm1\}\). Then the \(K_1\)-class
\(\partial_Y(E_+(\pi(P_\pm)))\) of Theorem~\ref{thm:pathb_roe_winding} maps under the
Toeplitz extension to the integer
\[
\operatorname{wind}\bigl(\partial_Y(E_+(\pi(P_\pm)))\bigr)
=\mp\,\varepsilon(t,n)\Ch\bigl(E_-(H_\pm^\triangle)\bigr)=\pm1,
\]
identifying the \(K_1\)-class with \(\mathbb Z\) via the Toeplitz quotient in
Theorem~\ref{thm:periodic_toeplitz_pairing}. The integer-pairing statement in this case is
fully internal: the closed-bulk gap and Chern signs are proved in
Theorem~\ref{thm:periodic_triangular_bulk} and Theorem~\ref{thm:periodic_triangular_chern},
and the Toeplitz identification of \(K_1\) with \(\mathbb Z\) is the standard BDF--Toeplitz
extension recalled in Theorem~\ref{thm:periodic_toeplitz_pairing} and
Corollary~\ref{cor:primitive_periodic_toeplitz}.
\end{corollary}
\proofsuppref{app:proof:cor:toeplitz_winding}

\begin{corollary}[Conditional integer winding for general partitions]\label{cor:conditional_winding}
Suppose that the \(K_1\)-class of the interface ideal \(\mathcal I=C^*(Y\subset X)\) admits
an integer-valued pairing \(\langle\,\cdot\,,\,\cdot\,\rangle:K_0(\mathcal A/\mathcal I)\times
K_1(\mathcal I)\to\mathbb Z\) of Chern-jump type, in the sense that
\(\langle [p_J],\partial_Y([p_J])\rangle=\Ch([P_+])-\Ch([P_-])\) whenever both sides are
defined and \(\partial_Y\) is the Higson--Roe connecting map of
Convention~\ref{conv:pathb_sign}. Then the pairing of
\(\partial_Y(E_+(\pi(P_\pm)))\) with \([p_J]\) is \(\pm1\). The Toeplitz periodic case
(Corollary~\ref{cor:toeplitz_winding}) is the special case where the integer pairing is
internal; the general case is conditional on the existence of an Ewert--Meyer/Chern-jump
compatible integer pairing whose construction outside the periodic class is not carried out in
this paper.
\end{corollary}
\proofsuppref{app:proof:cor:conditional_winding}

\begin{corollary}[Finite-\(M\) cutoff stability]\label{cor:pathb_finite_m_winding}
For every \(M>B_{\mathrm{copy},-}\), the \(K_1(\mathcal I)\)-class
\(\partial_Y(E_+(\pi(P_\pm)))\) of Theorem~\ref{thm:pathb_roe_winding} is independent of the
jam scale \(M\): equality of classes in \(K_1(\mathcal I)\) follows from the cutoff
homotopy and jam-scale-independence statements
(Corollary~\ref{S-cor:jam_independence} and Proposition~\ref{S-prop:quotient_sidewise}).
\end{corollary}

\subsection{Stability and cross-checks}

\begin{lemma}[Collar-supported stability without norm-smallness]\label{lem:pathb_collar_stability}
Let $V$ be any bounded self-adjoint perturbation in $C^*(Y\subset X)$, with no small-norm
assumption, and suppose the positive and negative comparison gaps used to define the sidewise
interface remain open. Then the quotient projection and \(K_1\)-class
\(\partial_Y(E_+(\pi(P_\pm)))\) of Theorem~\ref{thm:pathb_roe_winding} are unchanged. In the
straight periodic Toeplitz case (Corollary~\ref{cor:toeplitz_winding}), the resulting integer
winding is also unchanged.
\end{lemma}
\proofsuppref{app:proof:lem:pathb_collar_stability}

\begin{lemma}[Bounded-geometry stability]\label{lem:pathb_bg_stability}
Along a norm-continuous bounded-geometry deformation of the closed doubled universal Hamiltonians
that preserves the Higson--Prodan Fermi gap, the interface \(K_1\)-class of
Theorem~\ref{thm:pathb_roe_winding} is constant.
\end{lemma}
\proofsuppref{app:proof:lem:pathb_bg_stability}

\begin{lemma}[Refinement stability]\label{lem:pathb_refinement_stability}
For a refinement family $K_n\to K_\infty$ in bounded geometry, if the Higson--Prodan theorem gives
a uniform gap $g_{\mathrm{HP}}$ and the jam scales satisfy $M_n>B_{\mathrm{copy},-}+1$,
$M_n\to\infty$, then the quotient class and \(K_1\)-class of
Theorem~\ref{thm:pathb_roe_winding} are eventually constant under subsequences or cofinal
refinements.
\end{lemma}
\proofsuppref{app:proof:lem:pathb_refinement_stability}

\begin{remark}[Periodic triangular cross-check]\label{lem:pathb_periodic_triangular}
For periodic triangular doubles, Theorem~\ref{thm:pathb_roe_winding} agrees with the internal
Bloch gap, Chern, and straight Toeplitz winding calculations of
Section~\ref{sec:periodicbulk}; details in the supplement,
Section~\ref{S-app:proof:lem:pathb_periodic_triangular}.
\end{remark}

\begin{remark}[Finite-defect cross-check]\label{lem:pathb_finite_defect}
For doubled finite-defect families whose defects remain in a fixed collar or in finitely many
coarse components, Theorem~\ref{thm:pathb_roe_winding} gives the same interface class
as the finite-defect robustness statements in the supplement,
\Cref{S-app:robustness_finite_defects}; details in
Section~\ref{S-app:proof:lem:pathb_finite_defect}.
\end{remark}

\begin{remark}[Hexagonal and square cross-check]\label{lem:pathb_hex_square}
For the supplement's hexagonal and square symbols (\Cref{S-app:benchmark_hex_square}), whenever
their closed-bulk class is identified with the corresponding Higson--Prodan closed-bulk class or
with the internally computed periodic Chern class through a gapped homotopy,
Theorem~\ref{thm:pathb_roe_winding} returns the same interface \(K_1\)-class, and the
Toeplitz integer winding agrees with Corollary~\ref{cor:toeplitz_winding} in the periodic
case; details in Section~\ref{S-app:proof:lem:pathb_hex_square}.
\end{remark}

\begin{remark}[Tail-split and annular cross-check]\label{lem:pathb_tail_annular}
For the tail-split and annular families of the supplement
(\Cref{S-app:robustness_tail_annular}), the boundary class of
Theorem~\ref{thm:pathb_roe_winding} restricts componentwise to the same Toeplitz quotient
coordinates used there; details in Section~\ref{S-app:proof:lem:pathb_tail_annular}.
\end{remark}

\begin{remark}[Input synthesis]\label{rem:pathb_synthesis}
A naive accounting separates the closed doubled-bulk gap, the closed-bulk Chern pairing, and the
Ewert--Meyer comparison into three distinct inputs. Theorem~\ref{thm:pathb_roe_winding} packages
the first two as the single Higson--Prodan closed theorem and proves the third for the
doubled-jammed original/copy partition. The only cited closed-model theorem needed for the
\(K_1\)-class identity is the Higson--Prodan theorem. The integer winding evaluation is
internal in the periodic triangular Toeplitz case (\Cref{cor:toeplitz_winding}) and is
recorded as a conditional corollary for general partitions
(\Cref{cor:conditional_winding}).
\end{remark}

The companion coarse interface class construction, which places the doubled-and-jammed refinement
family in a Roe-algebra construction producing the same quotient projection and \(K_1\)-class
from uniform locality, boundedness, and an available closed-bulk gap, is recorded in the
supplement, Section~\ref{S-sec:bulkboundary}. Its main statement
(Theorem~\ref{S-thm:bulkboundary_main}) and the Ewert--Meyer-language restatement
(Theorem~\ref{S-thm:full_hp_bec}) consolidate the same conclusion.

%====================================================================
\section{A periodic triangular bulk input}\label{sec:periodicbulk}
%====================================================================

One closed-bulk input is checked directly here, rather than imported from Higson--Prodan. Let
\(z_j=e^{ik_j}\), \(k=(k_1,k_2)\in\mathbb T^2\), and consider the
\(\mathbb Z^2\)-periodic triangular triangulation with one vertex \(v\), three oriented edges
\(a,b,c\), and two oriented faces \(f,g\) per fundamental cell. The Bloch fiber is
\[
\mathcal H(k)
=
\mathbb C\langle v\rangle\oplus
\mathbb C\langle a,b,c\rangle\oplus
\mathbb C\langle f,g\rangle .
\]
With respect to this ordered basis, the boundary blocks are
\[
\partial_1(k)=
\begin{bmatrix}
z_1-1&z_2-1&z_1z_2-1
\end{bmatrix},
\qquad
\partial_2(k)=
\begin{bmatrix}
1&-z_2\\
z_1&-1\\
-1&1
\end{bmatrix}.
\]
A direct multiplication gives \(\partial_1(k)\partial_2(k)=0\). The cap-product blocks are
\[
P_0(k)=
\begin{bmatrix}
z_1^{-1}z_2^{-1}\\
z_2^{-1}
\end{bmatrix},
\qquad
P_1(k)=
\begin{bmatrix}
0&z_1^{-1}&0\\
0&0&0\\
-z_2^{-1}&0&0
\end{bmatrix},
\qquad
P_2(k)=
\begin{bmatrix}
1&1
\end{bmatrix},
\]
and the Hermitian duality block is
\[
S_0(k)=\tfrac{i}{2}(P_2(k)^*+P_0(k)),\quad
S_1(k)=\tfrac{i}{2}(P_1(k)^*-P_1(k)),\quad
S_2(k)=-\tfrac{i}{2}(P_0(k)^*+P_2(k)),
\]
so that \(H_\pm^\triangle(k)=B(k)+B(k)^*\pm S(k)\) is Hermitian.

\begin{proposition}[Periodic triangular characteristic polynomial]\label{prop:periodic_triangular_charpoly}
Set
\[
F(k)=28-8\cos k_1-6\cos k_2-6\cos(k_1+k_2),
\qquad
G(k)=\sin k_2+\sin(k_1+k_2).
\]
Then
\[
\det\!\left(\lambda\Id-H_\pm^\triangle(k)\right)
=
\bigl(\lambda^2-\tfrac{F(k)}4\bigr)
\bigl(\lambda^4-\tfrac{13}{2}\lambda^2+\tfrac{F(k)}4
\pm \lambda G(k)\bigr),
\]
with one exact pair of eigenvalue branches \(\lambda=\pm \tfrac12\sqrt{F(k)}\).
\end{proposition}
\proofsuppref{app:proof:prop:periodic_triangular_charpoly}

\begin{theorem}[Periodic triangular bulk gap]\label{thm:periodic_triangular_bulk}
For every \(k\in\mathbb T^2\), \(H_\pm^\triangle(k)\) is invertible. More precisely,
\[
\det H_\pm^\triangle(k)
=
-\tfrac1{16}
\left(
-28+8\cos k_1+6\cos k_2+6\cos(k_1+k_2)
\right)^2
\le -4.
\]
Moreover \(\|H_\pm^\triangle(k)\|\le8\), hence
\(\dist\bigl(0,\spec H_\pm^\triangle(k)\bigr)\ge 2^{-13}\) for all \(k\in\mathbb T^2\), and the
negative Bloch bundle has rank \(3\).
\end{theorem}
\proofsuppref{app:proof:thm:periodic_triangular_bulk}

\begin{corollary}[Small periodic perturbations preserve the triangular gap]\label{cor:periodic_small_perturbation_gap}
If \(V(k)\) is a continuous self-adjoint \(6\times6\) matrix over \(\mathbb T^2\) with
\(\sup_k\opnorm{V(k)}<2^{-13}\), then \(H_\pm^\triangle(k)+V(k)\) is invertible for every
\(k\in\mathbb T^2\), and its negative Bloch bundle has rank \(3\).
\end{corollary}
\proofsuppref{app:proof:cor:periodic_small_perturbation_gap}

\begin{proposition}[Positive duality strength homotopy]\label{prop:periodic_positive_mass_homotopy}
For \(m>0\), the operator \(H_{\pm,m}^\triangle(k)=B(k)+B(k)^*\pm mS(k)\) is invertible for every
\(k\in\mathbb T^2\); explicitly,
\[
\det H_{\pm,m}^\triangle(k)
=
-\tfrac1{16}
\bigl(8(3-c_1-c_2-c_{12})+2m^2(2+c_2+c_{12})\bigr)^2,
\]
with \(c_j=\cos k_j\), \(c_{12}=\cos(k_1+k_2)\). All positive values of the duality strength are
connected by a zero-gapped norm-continuous homotopy.
\end{proposition}
\proofsuppref{app:proof:prop:periodic_positive_mass_homotopy}

\begin{theorem}[Periodic triangular Chern signs]\label{thm:periodic_triangular_chern}
For the periodic triangular Bloch Hamiltonians above, the rank-three negative Bloch bundles have
\[
\Ch\bigl(E_-(H_+^\triangle)\bigr)=-1,
\qquad
\Ch\bigl(E_-(H_-^\triangle)\bigr)=+1.
\]
\end{theorem}
\proofsuppref{app:proof:thm:periodic_triangular_chern}

\begin{corollary}[Periodic Chern stability]\label{cor:periodic_chern_stability}
With \(g_{\pm,m}:=\min_k\dist(0,\spec H_{\pm,m}^\triangle(k))>0\), every continuous self-adjoint
periodic perturbation \(V(k)\) with \(\sup_k\opnorm{V(k)}<g_{\pm,m}\) keeps
\(H_{\pm,m}^\triangle(k)+V(k)\) gapped at zero and preserves the Chern number. The same holds for
\(H_\pm^\triangle\) when \(\sup_k\opnorm{V(k)}<2^{-13}\). In particular, sufficiently small
finite-range self-adjoint periodic perturbations of the triangular model do not change the Chern
signs.
\end{corollary}
\proofsuppref{app:proof:cor:periodic_chern_stability}

As a numerical check, the companion script applies the gauge-invariant lattice formula of
Fukui--Hatsugai--Suzuki \cite{FukuiHatsugaiSuzuki} to the rank-three negative spectral projection,
checking the returned integer, the link denominators, and the plaquette branch margins; the
computation records the same signs as Theorem~\ref{thm:periodic_triangular_chern}. The
antiunitary equivalence
\(H_-^\triangle(k)=\overline{H_+^\triangle(-k)}\) explains the opposite signs through complex
conjugation of the Bloch bundle, with the base map \(k\mapsto-k\) orientation-preserving in two
dimensions.

\begin{theorem}[Toeplitz evaluation for a straight periodic interface]\label{thm:periodic_toeplitz_pairing}
Let \(H(k_1,k_2)\) be a continuous finite-range self-adjoint \(N\times N\) Bloch Hamiltonian on
\(\mathbb T^2\), gapped at zero, and let
\(P(k_1,k_2)=\chi_{(-\infty,0)}(H(k_1,k_2))\). Write
\(\Ch(P)=\frac{1}{2\pi i}\int_{\mathbb T^2}\tr(P\,dP\wedge dP)\in\mathbb Z\) with the
orientation \(dk_1\wedge dk_2\). Compress the second lattice direction to a half-space. After
Fourier transform in the first direction, the straight-edge extension is the Toeplitz extension
\(0\to C(S^1,\mathcal K)\to\mathcal T_{\mathrm{per}}\xrightarrow{\sigma}C(\mathbb T^2)\to 0\) with
connecting map \(\partial_{\mathrm T}:K_0(C(\mathbb T^2))\to K_1(C(S^1))\) satisfying
\(\operatorname{wind}(\partial_{\mathrm T}([P]))=\Ch(P)\). The sign is fixed by the above
orientation and by taking the compressed coordinate to be the positive normal direction; see, for
example, \cite{BDF1977,BlackadarKTheory,ProdanSchulzBaldesBook,BourneKellendonkRennie}.
\end{theorem}
\proofsuppref{app:proof:thm:periodic_toeplitz_pairing}

\begin{corollary}[Primitive-direction straight periodic interfaces]\label{cor:primitive_periodic_toeplitz}
Let \(t,n\in\mathbb Z^2\) be a lattice basis with \(t\) the tangential direction of a straight
periodic interface and \(n\) the compressed positive normal direction. With
\(\varepsilon(t,n):=\det[t\ n]\in\{\pm1\}\), the half-space Toeplitz boundary map in the
\(\eta\)-direction (\(k=\kappa t+\eta n\)) satisfies
\(\operatorname{wind}\!\left(\partial_{\mathrm T,t,n}([P])\right)=\varepsilon(t,n)\,\Ch(P)\). In
particular, every primitive straight periodic edge direction is covered by a lattice change of
variables.
\end{corollary}
\proofsuppref{app:proof:cor:primitive_periodic_toeplitz}

\begin{corollary}[Periodic straight-edge winding from Toeplitz pairing]\label{cor:periodic_straight_edge_winding}
For the periodic triangular bulk models \(H^\triangle_\pm\), every primitive straight periodic
interface with tangential/positive-normal basis \((t,n)\) has interface winding
\[
\operatorname{wind}=-\det[t\ n]\quad\text{for }H^\triangle_+,
\qquad
\operatorname{wind}=\det[t\ n]\quad\text{for }H^\triangle_-.
\]
The same formula holds for every fixed positive-duality-strength triangular model and for every
small finite-range periodic perturbation covered by
Corollary~\ref{cor:periodic_chern_stability}. Finite-cover, finite-defect, finite-corner, split,
and tail-split variants are the appendix statements collected in the supplement,
Section~\ref{S-app:section5_variants}.
\end{corollary}
\proofsuppref{app:proof:cor:periodic_straight_edge_winding}

Hexagonal and square two-band benchmarks are recorded in the supplement,
Section~\ref{S-app:benchmark_hex_square}. Sign-ledger and orientation conventions are recorded in
the supplement, Section~\ref{S-app:sign_ledger}.

%====================================================================
\section{Numerical checks}\label{sec:numerical_diagnostics}
%====================================================================

The numerical material is collected in the online supplement, Section~\ref{S-sec:numerics}, with
reproducibility details in \Cref{S-app:reproducibility_appendix,S-app:minimal_reproduction_script}.
The computations are not proof inputs; they document finite-matrix behavior for the mechanisms
proved above. We retain two numerical figures in the main paper that bear directly on the
headline results: the intrinsic Poincar\'e--Lefschetz comparison
(Figure~\ref{fig:pl_boundary_comparison}) and the finite-\(M\) decoupling log-log scaling
(Figure~\ref{fig:finiteM_loglog}).

\medskip\noindent\textbf{Intrinsic boundary comparison.}
Figure~\ref{fig:pl_boundary_comparison} compares the intrinsic Poincar\'e--Lefschetz boundary
Hamiltonian \(H^{\mathrm{PL}}_+(\K)\), the doubled original-side compression \(H^\partial_+(\K)\),
and the finite-\(M\) jammed model on a disk-like sample. The collar correction
\(R_{\partial,+}=H^\partial_+(\K)-H^{\mathrm{PL}}_+(\K)\) is not norm-small, but it is
confined to the boundary collar predicted by
Theorem~\ref{thm:relative_boundary_comparison}, and the jammed low-energy spectrum tracks the
compressed boundary model in line with the Schur-complement bounds of
Theorem~\ref{thm:resolvent} and Theorem~\ref{thm:first_order_eff}. The numerical refinement scaling of
\(R_{\partial,+}\) -- recorded in the supplement -- is consistent with the boundary-rank
interpretation in Lemma~\ref{lem:collar_rank_counts}.

\begin{figure}[H]
\centering
\safeincludegraphics[width=0.92\linewidth]{fig_pl_boundary_comparison.png}
\caption{Central spectra of \(H^{\mathrm{PL}}_+(\K)\), \(H^\partial_+(\K)\), and
\(H_{+,M}\) on the disk-like example, together with the measured support distribution of the
collar correction \(R_{\partial,+}\). The correction is not norm-small, but it is boundary-collar
supported, as predicted by Theorem~\ref{thm:relative_boundary_comparison}.}
\label{fig:pl_boundary_comparison}
\end{figure}

\medskip\noindent\textbf{Finite-\(M\) log-log scalings.}
Figure~\ref{fig:finiteM_loglog} reports the finite-\(M\) rate diagnostics. The compressed
resolvent error follows the \(M^{-1}\) bound of Theorem~\ref{thm:resolvent}; the first-order
resolvent correction, squared copy leakage, and first-order eigenvalue localization follow the
\(M^{-2}\) predictions of \Cref{thm:first_order_eff,thm:eigvec_leak} and the first-order
eigenvalue localization theorem stated in the supplement,
\Cref{S-app:proof:thm:first_order_spectral_localization}.

\begin{figure}[H]
\centering
\begin{subfigure}{0.49\linewidth}
  \centering
  \safeincludegraphics[width=\linewidth]{fig_resolvent_loglog.png}
  \caption{Compressed resolvent.}
\end{subfigure}\hfill
\begin{subfigure}{0.49\linewidth}
  \centering
  \safeincludegraphics[width=\linewidth]{fig_first_order_resolvent_loglog.png}
  \caption{First-order resolvent.}
\end{subfigure}

\vspace{0.5em}
\begin{subfigure}{0.49\linewidth}
  \centering
  \safeincludegraphics[width=\linewidth]{fig_leakage_loglog.png}
  \caption{Copy leakage.}
\end{subfigure}\hfill
\begin{subfigure}{0.49\linewidth}
  \centering
  \safeincludegraphics[width=\linewidth]{fig_first_order_localization_loglog.png}
  \caption{First-order eigenvalue localization.}
\end{subfigure}
\caption{Log-log finite-\(M\) checks. The compressed resolvent slope tracks the \(M^{-1}\) rate
of Theorem~\ref{thm:resolvent}; the first-order resolvent, squared copy leakage, and first-order
eigenvalue localization track the \(M^{-2}\) rates of \Cref{thm:first_order_eff,thm:eigvec_leak}
and the first-order eigenvalue localization theorem stated in the supplement
(\Cref{S-app:proof:thm:first_order_spectral_localization}). Fitted slopes and additional rate
diagnostics are reported in the supplement, Section~\ref{S-sec:numerics}.}
\label{fig:finiteM_loglog}
\end{figure}

The remaining numerical checks -- boundary transport heatmaps, chirality flip, multiboundary
component drifts, flux insertion, collar/contour/refinement diagnostics, weak-disorder
boundary transport, and the nonconstant positive copy-side confinement -- are all collected in
the supplement, Section~\ref{S-sec:numerics}. A symbolic verification appendix for the local
crossing calculation supporting Theorem~\ref{thm:periodic_triangular_chern} is given in the
supplement, Section~\ref{S-app:symbolic_verification}.

\section{Conclusion}

We constructed a boundary model for the universal Chern Hamiltonians by replacing a triangulated
surface with boundary by its closed double and adding a large positive potential on the reflected
side. The relative fundamental class of \((\K,\bd\K)\) gave an intrinsic Poincar\'e--Lefschetz
boundary Hamiltonian on the original triangulation. Compressing the doubled Hamiltonian to the
original side recovered this relative Hamiltonian up to a boundary-collar operator, and the
jammed Hamiltonian converged to the compressed boundary operator in the large-jam limit by
quantitative Schur-complement estimates. The finite-\(M\) estimates controlled resolvents,
first-order self-energy corrections, spectral localization, leakage, contour projections,
Morse-index stability, monotone jam-scale spectral counts, no-copy low-energy bounds, and
inherited gaps under checkable separation conditions.

In the refinement limit, the Roe-algebra quotient by the interface ideal identified the
interface class as the connecting-map image of the positive-side doubled-bulk Fermi projection.
For sufficiently large \(M\), the jammed side contributed no negative-energy quotient sector.
Thus, for any boundary family whose closed doubles are covered by Higson--Prodan's published
theorem, the original/copy theorem of Section~\ref{sec:pathb_internalization} identifies the
interface \(K_1\)-class as \(\partial_Y(E_+(\pi(P_\pm)))\). In the periodic triangular Toeplitz
case this class evaluates internally to the integer winding \(\pm1\)
(Corollary~\ref{cor:toeplitz_winding}); for general partitions the integer pairing is recorded
as a conditional corollary (Corollary~\ref{cor:conditional_winding}). HP supplies the
closed-bulk gap and Chern value; the compatibility with the Roe boundary map is proved here. The periodic triangular section
gives a separate, fully internal bulk input: its Bloch symbol has an exact determinant gap, a
positive-mass homotopy remains gapped, and an isolated-crossing calculation proves the Chern
signs. For annular triangular-lattice families whose doubles are these periodic tori, the
boundary/interface construction therefore does not rely on the Higson--Prodan closed-bulk theorem
for the gap or Chern input.

Beyond the HP-covered and periodic classes treated here, one possible extension is to prove
uniform Fermi gaps for additional closed doubled refinement families. A second direction is to
extend the internal boundary-map identification to curved or componentwise pairings that separate
the contributions of individual boundary components.

\section*{Code and data availability}\label{sec:code_availability}
The companion numerical code and the online supplement are available at the project repository,
which is hosted on GitHub:
\href{https://github.com/yspennstate/Doubling-and-Jamming-A-Universal-Chern-Hamiltonian-on-Triangulated-Surfaces-with-Boundary}{\texttt{https://github.com/yspennstate/Doubling-and-Jamming-A-Universal-Chern-Hamiltonian-on-Triangulated-Surfaces-with-Boundary}}.
The online supplement is also distributed through that same GitHub repository. The
periodic triangular Bloch certificate script
(\texttt{periodic\_triangular\_bulk\_certificate.py}) and the companion numerical pipeline
that generates the figures and tables of the supplement are included.

\section*{Acknowledgments}
The research presented in this paper was supported by the European Research Council (ERC) under the European Union's Horizon Europe research and innovation programme (grant agreement No.~101041711), by the Simons Foundation (as part of the Collaboration on the Mathematical and Scientific Foundations of Deep Learning), by the Israel Science Foundation (grant number 2258/19), by the Israel Science Foundation (ISF Grant~4101/25), and by the U.S. National Science Foundation (NSF Grant~OISE-2401227).

\section*{Declaration of generative AI and AI-assisted technologies in the manuscript preparation process}
During the preparation of this work the author used generative AI tools to accelerate drafting
and revision. All mathematical claims, proofs, and bibliographic information were subsequently
reviewed and edited by the author, who takes full responsibility for the content of the manuscript.

\FloatBarrier

\begin{thebibliography}{99}

\bibitem{AsbothOroszlanyPalyi2016}
J.~K.~Asboth, L.~Oroszlany, and A.~Palyi,
\emph{A Short Course on Topological Insulators: Band Structure and Edge States in One and Two
Dimensions},
Lecture Notes in Physics \textbf{919}, Springer, 2016.

\bibitem{BellissardVanElstSchulzBaldes1994}
J.~Bellissard, A.~van Elst, and H.~Schulz-Baldes,
\emph{The noncommutative geometry of the quantum Hall effect},
J. Math. Phys. \textbf{35} (1994), 5373--5451.

\bibitem{Haldane1988}
F.~D.~M.~Haldane,
\emph{Model for a Quantum Hall Effect without Landau Levels: Condensed-Matter Realization of the Parity Anomaly},
Phys. Rev. Lett. \textbf{61} (1988), 2015--2018.

\bibitem{SchnyderRyuFurusakiLudwig}
A.~P.~Schnyder, S.~Ryu, A.~Furusaki, and A.~W.~W.~Ludwig,
\emph{Classification of topological insulators and superconductors in three spatial dimensions},
Phys. Rev. B \textbf{78} (2008), 195125.

\bibitem{KitaevPeriodic}
A.~Kitaev,
\emph{Periodic table for topological insulators and superconductors},
AIP Conf. Proc. \textbf{1134} (2009), 22--30.

\bibitem{QiZhangRMP}
X.-L.~Qi and S.-C.~Zhang,
\emph{Topological insulators and superconductors},
Rev. Mod. Phys. \textbf{83} (2011), 1057--1110.

\bibitem{FukuiHatsugaiSuzuki}
T.~Fukui, Y.~Hatsugai, and H.~Suzuki,
\emph{Chern numbers in discretized Brillouin zone: efficient method of computing (spin) Hall conductances},
J. Phys. Soc. Jpn. \textbf{74} (2005), 1674--1677.

\bibitem{HigsonProdanUniversalPRL}
N.~Higson and E.~Prodan,
\emph{Universal Chern Model on Arbitrary Triangulations},
Phys. Rev. Lett. \textbf{136} (2026), no.~9, 096602,
doi:10.1103/66x5-dvqq.

\bibitem{HigsonProdanUniversalArxiv}
N.~Higson and E.~Prodan,
\emph{A Universal Chern Model on Arbitrary Triangulations},
arXiv:2510.20862, doi:10.48550/arXiv.2510.20862.

\bibitem{HigsonRoeI}
N.~Higson and J.~Roe,
\emph{Mapping Surgery to Analysis I: Analytic Signatures},
K-Theory \textbf{33} (2004), 277--299.

\bibitem{HigsonRoeII}
N.~Higson and J.~Roe,
\emph{Mapping Surgery to Analysis II: Geometric Signatures},
K-Theory \textbf{33} (2004), 301--324.

\bibitem{HigsonRoeIII}
N.~Higson and J.~Roe,
\emph{Mapping Surgery to Analysis III: Exact Sequences},
K-Theory \textbf{33} (2004), 325--346.

\bibitem{WallSurgery}
C.~T.~C.~Wall,
\emph{Surgery on Compact Manifolds},
AMS, 1999.

\bibitem{MunkresElements}
J.~R.~Munkres,
\emph{Elements of Algebraic Topology},
Addison-Wesley, 1984.

\bibitem{Hatcher}
A.~Hatcher,
\emph{Algebraic Topology},
Cambridge University Press, 2002.

\bibitem{KatoPerturbation}
T.~Kato,
\emph{Perturbation Theory for Linear Operators},
Classics in Mathematics, Springer, 1995.

\bibitem{BhatiaMatrixAnalysis}
R.~Bhatia,
\emph{Matrix Analysis},
Graduate Texts in Mathematics, vol.~169, Springer, 1997.

\bibitem{AtiyahKTheory}
M.~F.~Atiyah,
\emph{K-Theory},
W.~A.~Benjamin, 1967.

\bibitem{BlackadarKTheory}
B.~Blackadar,
\emph{K-Theory for Operator Algebras},
2nd ed., Mathematical Sciences Research Institute Publications, vol.~5,
Cambridge University Press, 1998.

\bibitem{BDF1977}
L.~G.~Brown, R.~G.~Douglas, and P.~A.~Fillmore,
\emph{Extensions of $C^*$-algebras and $K$-homology},
Ann. of Math. (2) \textbf{105} (1977), no.~2, 265--324.

\bibitem{ProdanSchulzBaldesBook}
E.~Prodan and H.~Schulz-Baldes,
\emph{Bulk and Boundary Invariants for Complex Topological Insulators: From K-Theory to Physics},
Springer, 2016.

\bibitem{Mitchell2018}
N.~P.~Mitchell, L.~M.~Nash, D.~Hexner, A.~Turner, and W.~T.~M.~Irvine,
\emph{Amorphous topological insulators constructed from random point sets},
Nat. Phys. \textbf{14} (2018), 380--385.

\bibitem{BourneProdan2018}
C.~Bourne and E.~Prodan,
\emph{Non-commutative Chern numbers for generic aperiodic discrete systems},
J. Phys. A: Math. Theor. \textbf{51} (2018), 235202.

\bibitem{BourneMesland2019}
C.~Bourne and B.~Mesland,
\emph{Index theory and topological phases of aperiodic lattices},
Ann. Henri Poincar\'e \textbf{20} (2019), 1969--2038.

\bibitem{ElbauGraf2002}
P.~Elbau and G.~M.~Graf,
\emph{Equality of Bulk and Edge Hall Conductance Revisited},
Comm. Math. Phys. \textbf{229} (2002), 415--432.

\bibitem{ElgartGrafSchenker2005}
A.~Elgart, G.~M.~Graf, and J.~H.~Schenker,
\emph{Equality of the Bulk and Edge Hall Conductances in a Mobility Gap},
Comm. Math. Phys. \textbf{259} (2005), 185--221.

\bibitem{GrafPorta2013}
G.~M.~Graf and M.~Porta,
\emph{Bulk-Edge Correspondence for Two-Dimensional Topological Insulators},
Comm. Math. Phys. \textbf{324} (2013), 851--895.

\bibitem{KubotaControlled}
Y.~Kubota,
\emph{Controlled topological phases and bulk-edge correspondence},
Comm. Math. Phys. \textbf{349} (2017), 493--525.

\bibitem{EwertMeyer2019}
E.~E.~Ewert and R.~Meyer,
\emph{Coarse geometry and topological phases},
Comm. Math. Phys. \textbf{366} (2019), 1069--1098,
doi:10.1007/s00220-019-03303-z.

\bibitem{BourneKellendonkRennie}
C.~Bourne, J.~Kellendonk, and A.~Rennie,
\emph{The $K$-theoretic bulk-edge correspondence for topological insulators},
Ann. Henri Poincar\'e \textbf{18} (2017), 1833--1866.

\bibitem{AlldridgeMaxZirnbauer2020}
A.~Alldridge, C.~Max, and M.~R.~Zirnbauer,
\emph{Bulk-Boundary Correspondence for Disordered Free-Fermion Topological Phases},
Comm. Math. Phys. \textbf{377} (2020), 1761--1821.

\bibitem{LudewigThiangEdge}
M.~Ludewig and G.~C.~Thiang,
\emph{Cobordism invariance of topological edge-following states},
Adv. Theor. Math. Phys. \textbf{26} (2022), no.~3, 673--710.

\bibitem{DrouotZhuEdgeSpectrum}
A.~Drouot and X.~Zhu,
\emph{Topological edge spectrum along curved interfaces},
Int. Math. Res. Not. IMRN \textbf{2024} (2024), no.~22, 13870--13889.

\bibitem{DrouotZhuBEC}
A.~Drouot and X.~Zhu,
\emph{The bulk-edge correspondence for curved interfaces},
arXiv:2408.07950.

\bibitem{RoeLectures}
J.~Roe,
\emph{Lectures on Coarse Geometry},
University Lecture Series \textbf{31}, American Mathematical Society, 2003.

\end{thebibliography}

\end{document}
